% Generated draft: requires mathematical and source review.
\documentclass[11pt]{article}
\usepackage[margin=25mm]{geometry}
\usepackage{fontspec}
\setmainfont{Latin Modern Roman}
\usepackage{amsmath,amssymb,mathtools,unicode-math}
\setmathfont{Latin Modern Math}
\usepackage{enumitem,needspace,xurl,hyperref}
\setlength{\parindent}{0pt}
\setlength{\parskip}{5pt}
\setlength{\emergencystretch}{3em}
\newcommand{\N}{\mathbb{N}}
\newcommand{\Z}{\mathbb{Z}}
\newcommand{\Q}{\mathbb{Q}}
\newcommand{\R}{\mathbb{R}}
\newcommand{\C}{\mathbb{C}}
\newcommand{\F}{\mathbb{F}}
\newcommand{\A}{\mathbb{A}}
\newcommand{\PP}{\mathbb{P}}
\newcommand{\E}{\mathbb{E}}
\newcommand{\eps}{\varepsilon}
\newcommand{\dd}{\,\mathrm{d}}
\begin{document}
\section*{UnsolvedMath ID: 30001888}
\textbf{Source problem code:} catalogue entry
\subsection{Definitions and mathematical statement}
Let \(\R^2\) denote the Euclidean plane, and let \(D\subset\R^2\) be a disk. The supplied record asks whether
\[
\lambda_h(D,3)=\sqrt{3}
\]
for families with property \(T(3)\), describing \(\lambda_h(D,3)\) as an optimal homothety factor. The supplied background identifies the intended setting as a homothetic-copy variant.

\textbf{Unresolved formulation.} The record defines neither property \(T(3)\) nor the optimization quantity \(\lambda_h(D,3)\). It does not specify the admissible families, the geometric conclusion that the homothety factor must guarantee, the permitted centers and scale factors, or the quantifiers defining optimality. In particular, it does not establish whether \(T(3)\) concerns intersection or transversals. The title's reference to three intersecting disks does not resolve these omissions. The disk's boundary convention and whether degenerate disks are allowed are also unspecified. Thus the displayed equality records the proposed target, but a fully quantified mathematical statement cannot be recovered from the supplied material alone.
\subsection{Short English statement}
The proposed question is whether the optimal homothety factor for a disk and families satisfying \(T(3)\) equals \(\sqrt{3}\). Definitions of \(T(3)\), the admissible families, and the optimization problem are missing, so this remains an incomplete statement requiring mathematical review.
\subsection{Sources}
\begin{itemize}
\item Supplied local UnsolvedMath registry record, permanent ID \(30001888\), source problem code \texttt{catalogue entry}, titled ``Optimal Homothety Factor for Three Intersecting Disks,'' in the Oberwolfach Reports Open Problems collection. No individual proposer is supplied; the proposed year is recorded as 2011.
\item Supplied source citation: \emph{Discrete Geometry (2012)}, Oberwolfach Reports, DOI 10.4171/owr/2011/44, pp.\ 2459--2548. \url{https://doi.org/10.4171/owr/2011/44}. The background describes these reports as documenting mathematical workshops and their open-problem discussions. The report citation uses 2012, whereas the DOI and proposed-year field contain 2011; the supplied record does not explain this dating distinction or identify a problem-specific page.
\item Supplied literature source labeled ``OWR 44 (2012)'': \url{https://doi.org/10.4171/owr/2011/44}.
\item Supplied literature source labeled ``cover-decomposition survey'': \url{https://real.mtak.hu/15290/1/surveyfinal.pdf}.
\item Literature assessment dated exactly \texttt{2026-08-21}: the supplied curation labels the problem open, states that the cited translate-covering results do not supply the bounds sought for the homothetic-copy variant, and reports that no complete resolution was found. This is supplied dated curation, not an independently verified literature assessment. Its time-sensitive status and cited claims require independent verification.
\item The supplied statement status is \texttt{corrected\_verified}; its accompanying verification explains only that TeX notation was enclosed in mathematical delimiters without changing the mathematical content. It does not provide the missing definitions or establish mathematical correctness.
\end{itemize}
\end{document}
